\section{Regular Expressions}
\label{sec:regular-expressions}
\slidesrc{Syntax: RE \& word: 1--15}

Regular expressions are the first tool of the syntactic level
(Section~\ref{sec:course-map}). They are the simplest way to describe a \emph{set of
strings} by a pattern, and they are everywhere in NLP: tokenisers, sentence splitters,
data cleaning, rule-based chatbots, information extraction. Formally they describe the
\emph{regular languages}, the bottom level of the Chomsky hierarchy that we will meet in
Section~\ref{sec:grammars-parsing}.

This section covers:
\begin{enumerate}[nosep]
  \item the syntax of regular expressions: disjunction, negation, repetition, anchors
        (Section~\ref{sec:regex-syntax});
  \item a worked example that shows the two kinds of errors a pattern can make
        (Section~\ref{sec:regex-example});
  \item ELIZA: a whole ``conversation'' built from cascades of substitutions
        (Section~\ref{sec:eliza}).
\end{enumerate}

In these notes a pattern is written in colour, e.g.\ \re{[wW]oodchuck}, and the text it
matches is \underline{underlined}.

\subsection{The Syntax of Regular Expressions}
\label{sec:regex-syntax}
\slidesrc{Syntax: RE \& word: 3--8, 10}

\begin{nbdef}{Regular expression}
A \textbf{regular expression} (RE, regex) is a formal language for specifying text
strings. A regex is a \emph{pattern}; searching a text with it returns the substrings
that the pattern \emph{matches}.
\end{nbdef}

The motivating question: how can we search for any of \emph{woodchuck},
\emph{woodchucks}, \emph{Woodchuck}, \emph{Woodchucks} with a single query? A plain string
search for ``woodchuck'' misses the capitalised forms; we need a way to say ``either
\texttt{w} or \texttt{W}'' and ``an optional \texttt{s}''.

\subsubsection{Disjunction: square brackets and the pipe}
\slidesrc{Syntax: RE \& word: 4, 6}

Letters inside square brackets mean \textbf{one character among these}:

\begin{center}
\small
\begin{tabular}{@{}l l l@{}}
\toprule
\textbf{Pattern} & \textbf{Matches} & \textbf{Example} \\
\midrule
\re{[wW]oodchuck}   & \emph{Woodchuck} or \emph{woodchuck} & \\
\re{[1234567890]}   & any digit & \\
\re{[A-Z]}          & an upper-case letter (a \textbf{range}) & \underline{D}renched Blossoms \\
\re{[a-z]}          & a lower-case letter & \underline{m}y beans were impatient \\
\re{[0-9]}          & a single digit & Chapter \underline{1}: Down the Rabbit Hole \\
\bottomrule
\end{tabular}
\end{center}

Brackets choose among single \emph{characters}. To choose among longer strings use the
\textbf{pipe} \re{|}:

\begin{center}
\small
\begin{tabular}{@{}l l@{}}
\toprule
\textbf{Pattern} & \textbf{Matches} \\
\midrule
\re{groundhog|woodchuck} & \emph{groundhog} or \emph{woodchuck} (they are the same animal!) \\
\re{yours|mine} & \emph{yours} or \emph{mine} \\
\re{a|b|c} & the same as \re{[abc]} \\
\re{[gG]roundhog|[Ww]oodchuck} & both animals, capitalised or not \\
\bottomrule
\end{tabular}
\end{center}

\subsubsection{Negation inside brackets}
\slidesrc{Syntax: RE \& word: 5}

A caret \re{\textasciicircum} right after the opening bracket negates the set: \re{[\textasciicircum
Ss]} means ``any character that is \emph{neither} \texttt{S} nor \texttt{s}''.

\begin{center}
\small
\begin{tabular}{@{}l l l@{}}
\toprule
\textbf{Pattern} & \textbf{Matches} & \textbf{Example} \\
\midrule
\re{[\textasciicircum A-Z]} & not an upper-case letter & O\underline{y}fn pripetchik \\
\re{[\textasciicircum Ss]}  & neither \texttt{S} nor \texttt{s} & \underline{I} have no exquisite reason \\
\re{[\textasciicircum e\textasciicircum]} & neither \texttt{e} nor \texttt{\textasciicircum} & \underline{L}ook here \\
\re{a\textasciicircum b}    & the literal string \texttt{a\textasciicircum b} & Look up \underline{a\textasciicircum b} now \\
\bottomrule
\end{tabular}
\end{center}

\begin{nbwarning}[the caret has three meanings]
\re{\textasciicircum} means \emph{negation} only when it is the \emph{first} character inside
\re{[\,]}. In \re{[e\textasciicircum]} it is an ordinary character, in \re{a\textasciicircum b}
it is an ordinary character, and at the start of a pattern (\re{\textasciicircum[A-Z]}) it
is an \emph{anchor} (see below).
\end{nbwarning}

\subsubsection{Repetition: \texttt{?}, \texttt{*}, \texttt{+} and the wildcard \texttt{.}}
\slidesrc{Syntax: RE \& word: 7}

These operators apply to the character (or group) right before them:

\begin{center}
\small
\begin{tabular}{@{}l l l@{}}
\toprule
\textbf{Pattern} & \textbf{Meaning} & \textbf{Matches} \\
\midrule
\re{colou?r} & previous character optional (0 or 1) & \emph{color}, \emph{colour} \\
\re{oo*h!}   & 0 or more of the previous character & \emph{oh!}, \emph{ooh!}, \emph{oooh!}, \ldots \\
\re{o+h!}    & 1 or more of the previous character & \emph{oh!}, \emph{ooh!}, \emph{oooh!}, \ldots \\
\re{baa+}    & \texttt{ba} followed by one or more \texttt{a} & \emph{baa}, \emph{baaa}, \emph{baaaa}, \ldots \\
\re{beg.n}   & \texttt{.} is any single character & \emph{begin}, \emph{begun}, \emph{beg3n} \\
\bottomrule
\end{tabular}
\end{center}

The star \re{*} and the plus \re{+} are called the \textbf{Kleene star} and
\textbf{Kleene plus}, after the mathematician Stephen C.~Kleene who introduced regular
expressions. Note that \re{oo*h!} and \re{o+h!} describe exactly the same set of strings:
$x^{+}$ is just $x x^{*}$.

\subsubsection{Anchors: \texttt{\textasciicircum} and \texttt{\$}}
\slidesrc{Syntax: RE \& word: 8}

Anchors do not match a character but a \emph{position}: \re{\textasciicircum} the start
of the line, \re{\$} the end.

\begin{center}
\small
\begin{tabular}{@{}l l l@{}}
\toprule
\textbf{Pattern} & \textbf{Matches} & \textbf{Example} \\
\midrule
\re{\textasciicircum[A-Z]} & an upper-case letter at the start of a line & \underline{P}alo Alto \\
\re{\textasciicircum[\textasciicircum A-Za-z]} & a non-letter at the start of a line & \underline{1}\quad \underline{``}Hello'' \\
\re{\textbackslash.\$} & a full stop at the end of a line & The end\underline{.} \\
\re{.\$} & any character at the end of a line & The end\underline{?}\quad The end\underline{!} \\
\bottomrule
\end{tabular}
\end{center}

In \re{\textbackslash.\$} the backslash \textbf{escapes} the dot: \re{\textbackslash.}
is a literal full stop, while \re{.} alone is the wildcard.

\subsubsection{Summary of the operators}
\slidesrc{Syntax: RE \& word: 10}

\begin{center}
\small
\begin{tabularx}{\linewidth}{@{}l l l X@{}}
\toprule
\textbf{Operator} & \textbf{Meaning} & \textbf{Example} & \textbf{Example meaning} \\
\midrule
\re{+}     & one or more   & \re{a+}    & 1 or more \texttt{a} characters \\
\re{*}     & zero or more  & \re{a*}    & 0 or more \texttt{a} characters \\
\re{?}     & optional      & \re{a?}    & 0 or 1 \texttt{a} characters \\
\re{[\,]}  & choose 1      & \re{[abc]} & \texttt{a} or \texttt{b} or \texttt{c} \\
\re{[-]}   & range         & \re{[a-z]} & any character between \texttt{a} and \texttt{z} \\
\re{[\textasciicircum]} & not & \re{[\textasciicircum a]} & any character that is not \texttt{a} \\
\re{(\,)}  & grouping      & \re{([a-z])+} & one or more characters between \texttt{a} and \texttt{z} \\
\re{(|)}   & or            & \re{(ey|ax)} & the string \texttt{ey} or the string \texttt{ax} \\
\re{ab}    & follow        & \re{ab}    & \texttt{a} followed by \texttt{b} \\
\re{\textasciicircum} & start & \re{\textasciicircum a} & \texttt{a} at the start of the string/line \\
\re{\$}    & end           & \re{a\$}   & \texttt{a} at the end of the string/line \\
\re{\textbackslash s} & whitespace & \re{\textbackslash sa} & a whitespace followed by \texttt{a} \\
\re{.}     & any character & \re{a.b}   & \texttt{a}, any character, \texttt{b} \\
\bottomrule
\end{tabularx}
\end{center}

\begin{nbwarning}[a typo in the summary slide]
The slide gives \re{(a-z)+} as the grouping example. Inside round brackets the hyphen is
\emph{not} a range: \re{(a-z)+} matches the literal three-character string
\texttt{a-z} (repeated). The range needs square brackets: \re{([a-z])+} or simply
\re{[a-z]+}.
\end{nbwarning}

Tools to experiment: \url{https://regex101.com/} shows step by step what a pattern
matches and explains each part of it.

\subsection{Worked Example: Finding ``the''}
\label{sec:regex-example}
\slidesrc{Syntax: RE \& word: 9}

\textbf{Task}: find all instances of the word ``the'' in a text. Writing the right pattern
takes a few attempts, and each attempt fixes one kind of error. We test them on
\begin{center}
\emph{The other day the theology student said: the end.}
\end{center}

\begin{nbexample}{three attempts}
\begin{enumerate}
  \item \re{the}: matches \emph{the} three times, but also the \emph{the} inside
        \emph{o\underline{the}r} and \emph{\underline{the}ology}, and it
        \textbf{misses} the capitalised \emph{The} at the start.
  \item \re{[tT]he}: now \emph{The} is found, but it still \textbf{incorrectly returns}
        the \emph{the} inside \emph{other} and \emph{theology}.
  \item \re{[\textasciicircum a-zA-Z][tT]he[\textasciicircum a-zA-Z]}: ``the'' must be
        surrounded by non-letters, so \emph{other} and \emph{theology} are excluded.
\end{enumerate}
\end{nbexample}

Two kinds of errors appeared, and they are exactly the two denominators of precision and
recall (Section~\ref{sec:precision-recall}):
\begin{itemize}[nosep]
  \item \textbf{false positives}: matching strings we should not have matched
        (\emph{other}, \emph{theology}). Reducing them increases \textbf{precision};
  \item \textbf{false negatives}: not matching strings we should have matched
        (\emph{The}). Reducing them increases \textbf{recall}.
\end{itemize}
Fixing one kind of error often introduces the other, so building a pattern is a
continuous trade-off between the two.

\begin{nbwarning}[attempt 3 is still not perfect]
\re{[\textasciicircum a-zA-Z][tT]he[\textasciicircum a-zA-Z]} requires a character
\emph{before} and \emph{after} the word, so it misses ``The'' at the very beginning of a
line (there is no character before it) and a ``the'' at the very end. It also includes
the surrounding spaces in the match. Two common fixes:
\begin{itemize}[nosep]
  \item allow the start/end of the line as alternatives:
        \re{(\textasciicircum|[\textasciicircum a-zA-Z])[tT]he([\textasciicircum a-zA-Z]|\$)};
  \item use the \emph{word boundary} anchor \re{\textbackslash b} (available in Python
        and most regex engines): \re{\textbackslash b[tT]he\textbackslash b} matches
        \emph{The}, \emph{the}, \emph{the} and nothing else in the test sentence.
\end{itemize}
\end{nbwarning}

\subsection{ELIZA: Conversation by Substitution}
\label{sec:eliza}
\slidesrc{Syntax: RE \& word: 11--15}

ELIZA (Section~\ref{sec:nlp-history}) is the classic proof of how far regular
expressions alone can go. It understands nothing: it applies a \textbf{cascade} of
regex substitutions to the user's sentence and returns the result. A conversation with
an online version (\url{https://www.masswerk.at/elizabot/}) looks like this:

\begin{center}
\begin{minipage}{0.85\linewidth}
\small\ttfamily
ELIZA: How do you do. Please tell me your problem.\\
YOU:\ \ \ I had an unpleasant discussion with a friend of mine\\
ELIZA: Can you elaborate on that ?\\
YOU:\ \ \ She told me that I am stupid\\
ELIZA: Is it because you are stupid that you came to me ?\\
YOU:\ \ \ yes, it is exactly because I am stupid...\\
ELIZA: I see.
\end{minipage}
\end{center}

\emph{``Is it because you are stupid that you came to me?''} is just the user's
\emph{``I am stupid''} with \emph{I am} rewritten to \emph{you are} and pasted into a
template.

\subsubsection{A tiny ELIZA in Python}
\slidesrc{Syntax: RE \& word: 14--15}

Three substitutions applied in cascade with Python's \texttt{re.sub}, which takes a
pattern, a replacement and a string:

\begin{minted}{python}
import re
str = "I feel sad, my heart is broken"
str = re.sub("I", "you", str)        # you feel sad, my heart is broken
str = re.sub("my", "your", str)      # you feel sad, your heart is broken
str = re.sub(r"you (feel|are) (sad|sorry|desperate)(.*)",
             r"Why do you say that you \1 \2\3?", str)
print(str)   # Why do you say that you feel sad, your heart is broken?
\end{minted}

\begin{itemize}[nosep]
  \item Steps 1 and 2 swap the point of view (\emph{I} $\to$ \emph{you},
        \emph{my} $\to$ \emph{your}).
  \item Step 3 uses \textbf{capture groups}: each pair of round brackets in the pattern
        remembers what it matched, and \texttt{\textbackslash 1}, \texttt{\textbackslash 2},
        \texttt{\textbackslash 3} in the replacement paste it back. Here
        \texttt{\textbackslash 1} = \emph{feel}, \texttt{\textbackslash 2} = \emph{sad},
        \texttt{\textbackslash 3} = \emph{, your heart is broken}.
  \item The \texttt{r"..."} prefix makes a \emph{raw string}, so that Python does not
        interpret the backslashes itself.
\end{itemize}

\begin{nbwarning}[order and context matter]
\texttt{re.sub("I", "you", str)} replaces \emph{every} capital \texttt{I}, also inside
words (``\emph{It}'' becomes ``\emph{yout}''), and the order of the cascade matters: if
we swapped \emph{you} $\to$ \emph{I} as well, a later rule could undo an earlier one. Real
ELIZA implementations use word boundaries (\re{\textbackslash bI\textbackslash b}) and
swap pronouns in a single pass.
\end{nbwarning}

\subsubsection{Exercises}
\notinslides

\begin{nbexercise}{1}{prices}
Write a regex matching prices in dollars such as \emph{\$5} and \emph{\$19.99} (an
optional decimal part with exactly two digits). What does it return on ``not \$3.5''?
\end{nbexercise}
\begin{nbsolution}
\re{\textbackslash\$[0-9]+(\textbackslash.[0-9][0-9])?}: an escaped dollar sign (\re{\$}
alone is the end anchor), one or more digits, then an optional group made of a literal
dot and two digits. On ``\$3.5'' the optional group cannot match (only one digit after the
dot), so the regex returns just \emph{\$3}: a false positive of a wrong length.
\emph{Lesson}: a pattern can match only part of what you intended; add
\re{\textbackslash b} or a negative condition after the price if partial matches matter.
\end{nbsolution}

\begin{nbexercise}{2}{reading patterns}
Describe in words what each pattern matches: (a) \re{\textasciicircum[A-Z][a-z]*\textbackslash.\$};
(b) \re{[0-9]+(st|nd|rd|th)}; (c) \re{\textbackslash bcolou?rs?\textbackslash b}.
\end{nbexercise}
\begin{nbsolution}
(a) A whole line made of one capitalised word followed by a full stop (``Stop.'' but not
``Stop now.''). (b) An ordinal written with digits: \emph{1st}, \emph{22nd},
\emph{3rd}, \emph{100th} (and also the wrong \emph{1th}: a regex checks form, not
meaning). (c) The whole words \emph{color}, \emph{colour}, \emph{colors},
\emph{colours}, but not \emph{discolor}.
\end{nbsolution}

\begin{nbsummary}[regular expressions]
\begin{itemize}[nosep]
  \item \re{[abc]} one of; \re{[a-z]} range; \re{[\textasciicircum a]} not (caret first
        in brackets); \re{x|y} one of two strings; \re{(\,)} grouping and capture.
  \item \re{?} 0 or 1; \re{*} 0 or more (Kleene star); \re{+} 1 or more (Kleene plus);
        \re{.} any character; \re{\textbackslash.} a literal dot.
  \item Anchors \re{\textasciicircum} (start), \re{\$} (end), \re{\textbackslash b} (word
        boundary).
  \item Writing a pattern trades off \textbf{false positives} (lower precision) and
        \textbf{false negatives} (lower recall).
  \item ELIZA = a cascade of regex substitutions with capture groups
        (\texttt{\textbackslash 1}, \texttt{\textbackslash 2}, \ldots).
\end{itemize}
\end{nbsummary}

